\documentclass[11pt,paper=a4,answers]{exam}
\usepackage{graphicx,lastpage}
\usepackage{upgreek}
\usepackage{censor}
\usepackage{gensymb}
\usepackage{amsmath}
\usepackage{multicol}
\usepackage{enumitem}
\usepackage{tasks}
\usepackage{adjustbox}
\usepackage{xcolor}
\usepackage{tfrupee}
\setlength{\voffset}{0.25in}
\usepackage{tabularx}
\newcolumntype{Y}{>{\centering\arraybackslash}X}
%==============================================================
% Customise Non-Essential Packages Below
% =============================================================
\usepackage[most]{tcolorbox}
\usepackage{eso-pic}
\usepackage{tikz}
\AddToShipoutPictureBG{%
\begin{tikzpicture}[remember picture,overlay]
\foreach \x in {-8,-4,0,4,8,12,16,20}
\foreach \y in {-8,-4,0,4,8,12,16,20} {
\node[opacity=0.05,rotate=45,anchor=center] at ([xshift=\x cm,yshift=\y cm]current page.center) {%
\begin{minipage}{2cm}
\centering
\includegraphics[width=2cm]{images/logo_black.png}\\
\small preppeo.com
\end{minipage}%
};
}
\end{tikzpicture}%
}
\printanswers

%==============================================================
% Customise Non-Essential Packages Above
% =============================================================
\definecolor{svgray}{rgb}{0.90, 0.90, 0.90}
\graphicspath{{images/}}

\usepackage[normalem]{ulem}
\renewcommand\ULthickness{1pt}   %%---> For changing thickness of underline
\setlength\ULdepth{3.5ex}%\maxdimen ---> For changing depth of underline
\renewcommand{\baselinestretch}{1}
\chead{
\includegraphics[scale=0.1,height=2cm ]{logo.png}
}
\pagestyle{empty}

\pagestyle{headandfoot}
\headrule
\cfoot{W: https://preppeo.com; M: +91-8130711689}

\begin{document}
%==============================================================
% Customise Question paper details below this
% =============================================================
\begin{center}
    \centering
    \renewcommand{\arraystretch}{1.5}
    \begin{tabularx}{\textwidth}{YYY}
    & \normalsize\bf{{\Large{{Quantitative Reasoning}}}} & \\
    By Preppeo & \normalsize\bf{{sat\_lid\_034}} & SAT\\
    \normalsize{{\textbf{{Unit:}} One-Variable Data}} & \normalsize{{\textbf{{Chapter:}} Distributions, measures of center/spread}} & \normalsize{{\textbf{{Topic:}}Mean, Median, Mode, Range}} \\
    \end{tabularx}
\end{center}
\par\noindent\rule{\textwidth}{1pt}
% =============================================================
% Customise Question paper details above this
% =============================================================

\begin{questions}
\pointsinrightmargin
\pointsdroppedatright
\marksnotpoints
\pointpoints{mark}{marks}
\pointformat{\boldmath\themarginpoints}
\bracketedpoints

% =============================================================
% Question Begin
% =============================================================
% Question 1
% Category: Practice | Difficulty: Easy | Question Type: MCQ
\question
4, 10, 18, 4, 4, 5, 6, 5 \\
What is the median of the data set shown?
\begin{tasks}(4)
    \task 4
    \task 5
    \task 7
    \task 14
\end{tasks}

\begin{solution}
\textbf{Conceptual Explanation:} \\
The median is the middle value of an ordered data set. For an even number of values, it is the average of the two middle terms.

\vspace{30pt}

\textbf{Calculation and Logic:} \\
1. \textbf{Order the data:} 4, 4, 4, 5, 5, 6, 10, 18.
2. \textbf{Identify middle pair:} The 4th and 5th values are 5 and 5.
3. \textbf{Average:} $\frac{5+5}{2} = 5$.

\vspace{30pt}

Answer: \colorbox{yellow!100}{B}
\end{solution}

\vspace{60pt}

% Question 2
% Category: Practice | Difficulty: Easy | Question Type: MCQ
\question
2, 2, 2, 3, 4, 4, 11 \\
What is the median of the seven data values shown?
\begin{tasks}(4)
    \task 2
    \task 3
    \task 4
    \task 9
\end{tasks}

\begin{solution}
\textbf{Calculation and Logic:} \\
The data is ordered: 2, 2, 2, \textbf{3}, 4, 4, 11. The middle value (4th position) is 3.

\vspace{30pt}

Answer: \colorbox{yellow!100}{B}
\end{solution}

\vspace{60pt}

% Question 3
% Category: Exam Level | Difficulty: Medium | Question Type: MCQ
\question
The box plots summarize the masses, in kilograms, of two groups of gazelles. Based on the box plots, which of the following statements must be true?
\begin{tasks}(1)
    \task The mean mass of group 1 is greater than the mean mass of group 2.
    \task The mean mass of group 1 is less than the mean mass of group 2.
    \task The median mass of group 1 is greater than the median mass of group 2.
    \task The median mass of group 1 is less than the median mass of group 2.
\end{tasks}

\begin{solution}
\textbf{Calculation and Logic:} \\
1. Group 1 Median: 25 kg. \\
2. Group 2 Median: 24 kg. \\
3. Since $25 > 24$, the median of Group 1 is greater.

\vspace{30pt}

Answer: \colorbox{yellow!100}{C}
\end{solution}

\vspace{60pt}

% Question 4
% Category: Exam Level | Difficulty: Hard | Question Type: MCQ
\question
Ages of 20 Students Enrolled in a College Class
\begin{center}
\begin{tabular}{|c|c|}
\hline
Age & Frequency \\ \hline
18 & 6 \\ \hline
19 & 5 \\ \hline
20 & 4 \\ \hline
21 & 2 \\ \hline
22 & 1 \\ \hline
23 & 1 \\ \hline
30 & 1 \\ \hline
\end{tabular}
\end{center}
Which of the following gives the correct order of the mean, median, and mode of the ages?
\begin{tasks}(2)
    \task mode $<$ median $<$ mean
    \task mode $<$ mean $<$ median
    \task median $<$ mode $<$ mean
    \task mean $<$ mode $<$ median
\end{tasks}

\begin{solution}
\textbf{Calculation:} \\
1. \textbf{Mode:} 18 (highest frequency). \\
2. \textbf{Median:} 10th and 11th values are both 19. \\
3. \textbf{Mean:} $\frac{(18\cdot6)+(19\cdot5)+(20\cdot4)+(21\cdot2)+22+23+30}{20} = \frac{400}{20} = 20$. \\
Comparison: $18 < 19 < 20$.

\vspace{30pt}

Answer: \colorbox{yellow!100}{A}
\end{solution}

\vspace{60pt}

% Question 5
% Category: Practice | Difficulty: Medium | Question Type: MCQ
\question
Data set A: 72, 73, 73, 76, 76 \\
Data set B: 61, 64, 74, 85, $x$ \\
If the mean of data set A is equal to the mean of data set B, what is the value of $x$?
\begin{tasks}(4)
    \task 77
    \task 85
    \task 86
    \task 95
\end{tasks}

\begin{solution}
\textbf{Calculation:} \\
Sum A = 370. Sum B must be 370. \\
$61+64+74+85+x = 370 \Rightarrow 284+x=370 \Rightarrow x=86$.

\vspace{30pt}

Answer: \colorbox{yellow!100}{C}
\end{solution}

\vspace{60pt}

% Question 6
% Category: Exam Level | Difficulty: Medium | Question Type: MCQ
\question
Men's Height: Group A (Mean 186, SD 12.5), Group B (Mean 186, SD 19.1). Which statement is true?
\begin{tasks}(1)
    \task The Group A data set was identical to the Group B data set.
    \task Group B contained the tallest participant.
    \task Heights in Group B had a larger spread than in Group A.
    \task The median height of Group B is larger than Group A.
\end{tasks}

\begin{solution}
\textbf{Logic:} Standard deviation measures spread. Group B (19.1) $>$ Group A (12.5), so Group B has a larger spread.

\vspace{30pt}

Answer: \colorbox{yellow!100}{C}
\end{solution}

\vspace{60pt}

% Question 7
% Category: Exam Level | Difficulty: Hard | Question Type: MCQ
\question
Adding 7 to numbers above the median and subtracting 7 from numbers below the median. Which measure changes?
\begin{tasks}(2)
    \task Median
    \task Mean
    \task Sum of numbers
    \task Standard deviation
\end{tasks}

\begin{solution}
\textbf{Logic:} Values move further from the mean, increasing the spread. Standard deviation measures spread, so it will change.

\vspace{30pt}

Answer: \colorbox{yellow!100}{D}
\end{solution}

\vspace{60pt}

% Question 8
% Category: Practice | Difficulty: Medium | Question Type: Numeric Entry
\question
Class A range (0 to 5), Class B range (1 to 10). What is the positive difference between the ranges?

\begin{solution}
\textbf{Calculation:} Range A = 5. Range B = 9. Difference = $9 - 5 = 4$.

\vspace{30pt}

Answer: \colorbox{yellow!100}{4}
\end{solution}

\vspace{60pt}

% Question 9
% Category: Exam Level | Difficulty: Hard | Question Type: Numeric Entry
\question
Median of arrivals 2012 vs 2013 (9 values each). How much greater was the 2013 median?

\begin{solution}
\textbf{Calculation:} 2012 Median = 46.4. 2013 Median = 47.7. $47.7 - 46.4 = 1.3$.

\vspace{30pt}

Answer: \colorbox{yellow!100}{1.3}
\end{solution}

\vspace{60pt}

% Question 10
% Category: Exam Level | Difficulty: Medium | Question Type: MCQ
\question
[Dot Plots A and B] Which is true? I. Medians are equal. II. SDs are equal.
\begin{tasks}(4)
    \task I only
    \task II only
    \task I and II
    \task Neither
\end{tasks}

\begin{solution}
\textbf{Logic:} Both are symmetric at 13 (I is true). Set A has more dots at the extremes than B, so SDs are not equal (II is false).

\vspace{30pt}

Answer: \colorbox{yellow!100}{A}
\end{solution}

\vspace{60pt}

% Question 11
% Category: Concept | Difficulty: Medium | Question Type: MCQ
\question
A set of 5 positive integers has a mean of 10 and a median of 8. What is the largest possible value?
\begin{tasks}(4)
    \task 31
    \task 32
    \task 33
    \task 34
\end{tasks}

\begin{solution}
\textbf{Calculation:} Sum = 50. Set: 1, 1, 8, 8, $x$. $1+1+8+8+x=50 \Rightarrow x=32$.

\vspace{30pt}

Answer: \colorbox{yellow!100}{B}
\end{solution}

\vspace{60pt}

% Question 12
% Category: Exam Level | Difficulty: Hard | Question Type: MCQ
\question
Adding an outlier (10 siblings) to a set where max was 4. Which measure increases most?
\begin{tasks}(4)
    \task Mean
    \task Median
    \task Mode
    \task Range
\end{tasks}

\begin{solution}
\textbf{Logic:} Range changes from 4 to 10 (+6). Mean changes slightly. Median/Mode likely stay same.

\vspace{30pt}

Answer: \colorbox{yellow!100}{D}
\end{solution}

\vspace{60pt}

% Question 13
% Category: Practice | Difficulty: Medium | Question Type: Numeric Entry
\question
Set: 15, 18, 20, 22, 25. Add $k$ to make new mean 22. What is $k$?

\begin{solution}
\textbf{Calculation:} Old sum = 100. New sum ($22 \times 6$) = 132. $k = 132 - 100 = 32$.

\vspace{30pt}

Answer: \colorbox{yellow!100}{32}
\end{solution}

\vspace{60pt}

% Question 14
% Category: Concept | Difficulty: Hard | Question Type: MCQ
\question
Right-skewed distribution. Which is true?
\begin{tasks}(1)
    \task mean $<$ median $<$ mode
    \task mode $<$ median $<$ mean
    \task median $<$ mean $<$ mode
    \task mean = median = mode
\end{tasks}

\begin{solution}
\textbf{Logic:} Tail on right pulls mean highest. Mode is peak on left. Mode $<$ Median $<$ Mean.


\vspace{30pt}

Answer: \colorbox{yellow!100}{B}
\end{solution}

\vspace{60pt}

% Question 15
% Category: Exam Level | Difficulty: Medium | Question Type: MCQ
\question
Set X: 5, 10, 15, 20, 25. Set Y: 105, 110, 115, 120, 125. Compare SD.
\begin{tasks}(1)
    \task SD X is 100 less than SD Y
    \task SD X is equal to SD Y
    \task SD X is greater than SD Y
    \task SD X is 5 times less than SD Y
\end{tasks}

\begin{solution}
\textbf{Logic:} Adding a constant doesn't change spread. SD remains identical.

\vspace{30pt}

Answer: \colorbox{yellow!100}{B}
\end{solution}

\vspace{60pt}

% Question 16
% Category: Practice | Difficulty: Easy | Question Type: Numeric Entry
\question
Range of 10 numbers is 45. Smallest is 12. Largest?

\begin{solution}
\textbf{Calculation:} Max - 12 = 45 $\Rightarrow$ Max = 57.

\vspace{30pt}

Answer: \colorbox{yellow!100}{57}
\end{solution}

\vspace{60pt}

% Question 17
% Category: Exam Level | Difficulty: Hard | Question Type: MCQ
\question
Median is 20. Every number $\times 3$ then $+4$. New median?
\begin{tasks}(4)
    \task 20
    \task 60
    \task 64
    \task 75
\end{tasks}

\begin{solution}
\textbf{Calculation:} $(20 \times 3) + 4 = 64$.

\vspace{30pt}

Answer: \colorbox{yellow!100}{C}
\end{solution}

\vspace{60pt}

% Question 18
% Category: Concept | Difficulty: Medium | Question Type: MCQ
\question
Which is NOT affected by a single extreme outlier in a large set?
\begin{tasks}(4)
    \task Mean
    \task Range
    \task Standard Deviation
    \task Mode
\end{tasks}

\begin{solution}
\textbf{Logic:} Mode is the most frequent value; one outlier won't change frequency rankings in a large set.

\vspace{30pt}

Answer: \colorbox{yellow!100}{D}
\end{solution}

\vspace{60pt}

% Question 19
% Category: Exam Level | Difficulty: Medium | Question Type: MCQ
\question
Class A box is wider than Class B. Which is true?
\begin{tasks}(1)
    \task Class A has a higher median.
    \task IQR of Class A is greater.
    \task Class A has more students.
    \task Means are equal.
\end{tasks}

\begin{solution}
\textbf{Logic:} Box width = IQR ($Q3 - Q1$). A wider box means a larger IQR.


\vspace{30pt}

Answer: \colorbox{yellow!100}{B}
\end{solution}

\vspace{60pt}

% Question 20
% Category: Practice | Difficulty: Easy | Question Type: MCQ
\question
Mode of: 2, 5, 8, 2, 10, 5, 2, 7, 8
\begin{tasks}(4)
    \task 2
    \task 5
    \task 7
    \task 8
\end{tasks}

\begin{solution}
\textbf{Logic:} 2 appears three times; others appear less.

\vspace{30pt}

Answer: \colorbox{yellow!100}{A}
\end{solution}

\vspace{60pt}

% Question 21
% Category: Concept | Difficulty: Medium | Question Type: MCQ
\question
If the mean of a set is 50 and the SD is 0, which must be true?
\begin{tasks}(1)
    \task There is only one number in the set.
    \task All numbers in the set are 50.
    \task The median is 0.
    \task The range is 50.
\end{tasks}

\begin{solution}
\textbf{Logic:} SD = 0 means there is zero spread; every value must be identical to the mean.

\vspace{30pt}

Answer: \colorbox{yellow!100}{B}
\end{solution}

\vspace{60pt}

% Question 22
% Category: Exam Level | Difficulty: Medium | Question Type: MCQ
\question
A set of 7 numbers has a mean of 12. If 4 is added to the set, new mean?
\begin{tasks}(4)
    \task 8
    \task 10
    \task 11
    \task 12
\end{tasks}

\begin{solution}
\textbf{Calculation:} Old sum = $7 \times 12 = 84$. New sum = $84 + 4 = 88$. New $n = 8$. Mean = $88/8 = 11$.

\vspace{30pt}

Answer: \colorbox{yellow!100}{C}
\end{solution}

\vspace{60pt}

% Question 23
% Category: Practice | Difficulty: Hard | Question Type: Numeric Entry
\question
The average of $x, y,$ and $z$ is 15. The average of $x, y, z,$ and $w$ is 20. Value of $w$?

\begin{solution}
\textbf{Calculation:} Sum(3) = 45. Sum(4) = 80. $w = 80 - 45 = 35$.

\vspace{30pt}

Answer: \colorbox{yellow!100}{35}
\end{solution}

\vspace{60pt}

% Question 24
% Category: Concept | Difficulty: Medium | Question Type: MCQ
\question
Which measure is most sensitive to an outlier?
\begin{tasks}(4)
    \task Median
    \task IQR
    \task Range
    \task Mode
\end{tasks}

\begin{solution}
\textbf{Logic:} Range is determined solely by Max and Min. An extreme outlier immediately changes one of these.

\vspace{30pt}

Answer: \colorbox{yellow!100}{C}
\end{solution}

\vspace{60pt}

% Question 25
% Category: Exam Level | Difficulty: Hard | Question Type: MCQ
\question
Mean of 100 values is 70. If 10 values are removed (all being 70), new mean?
\begin{tasks}(4)
    \task 60
    \task 70
    \task 77
    \task 80
\end{tasks}

\begin{solution}
\textbf{Logic:} Removing values equal to the mean does not change the mean of the remaining set.

\vspace{30pt}

Answer: \colorbox{yellow!100}{B}
\end{solution}

% Question 26
% Category: Exam Level | Difficulty: Hard | Question Type: MCQ
\question
A data set of 150 values has a mean of 45 and a standard deviation of 8. If every value in the data set is increased by 10\%, what will be the new mean and the new standard deviation?
\begin{tasks}(1)
    \task Mean: 45, SD: 8
    \task Mean: 49.5, SD: 8
    \task Mean: 49.5, SD: 8.8
    \task Mean: 55, SD: 18
\end{tasks}

\begin{solution}
\textbf{Conceptual Explanation:} \\
When every value in a data set is multiplied by a constant $k$ (in this case, $1.10$), both the mean and the standard deviation are multiplied by that same constant.

\vspace{30pt}

\textbf{Calculation and Logic:} \\
1. \textbf{New Mean:} $45 \times 1.10 = 49.5$.
2. \textbf{New SD:} $8 \times 1.10 = 8.8$.

\vspace{30pt}

Answer: \colorbox{yellow!100}{C}
\end{solution}

\vspace{60pt}

% Question 27
% Category: Practice | Difficulty: Medium | Question Type: MCQ
\question
The median of a set of 9 distinct integers is 20. If the largest number is increased by 50, how does the median change?
\begin{tasks}(1)
    \task It increases.
    \task It decreases.
    \task It stays the same.
    \task It cannot be determined without the other values.
\end{tasks}

\begin{solution}
\textbf{Logic:} The median is the middle value. Changing the largest value (which is to the right of the median) does not shift the middle position or the value at that position.

\vspace{30pt}

Answer: \colorbox{yellow!100}{C}
\end{solution}

\vspace{60pt}

% Question 28
% Category: Exam Level | Difficulty: Medium | Question Type: MCQ
\question
Which of the following describes a data set where the mean is significantly greater than the median?
\begin{tasks}(1)
    \task The data is perfectly symmetric.
    \task The data is skewed to the left.
    \task The data is skewed to the right.
    \task The standard deviation is zero.
\end{tasks}

\begin{solution}
\textbf{Logic:} In a right-skewed distribution, high-value outliers pull the mean toward the right tail, while the median remains closer to the bulk of the data.

\vspace{30pt}

Answer: \colorbox{yellow!100}{C}
\end{solution}

\vspace{60pt}

% Question 29
% Category: Practice | Difficulty: Easy | Question Type: Numeric Entry
\question
What is the range of the following data set? \\
-12, -5, 0, 4, 15, 22

\begin{solution}
\textbf{Calculation:} Range = Max - Min = $22 - (-12) = 22 + 12 = 34$.

\vspace{30pt}

Answer: \colorbox{yellow!100}{34}
\end{solution}

\vspace{60pt}

% Question 30
% Category: Exam Level | Difficulty: Hard | Question Type: MCQ
\question
Two groups of students took a test. Group A ($n=20$) had a mean score of 80. Group B ($n=30$) had a mean score of 90. What is the mean score of all 50 students combined?
\begin{tasks}(4)
    \task 84
    \task 85
    \task 86
    \task 88
\end{tasks}

\begin{solution}
\textbf{Calculation:} \\
1. Total Score A = $20 \times 80 = 1600$.
2. Total Score B = $30 \times 90 = 2700$.
3. Total Combined = $1600 + 2700 = 4300$.
4. Combined Mean = $4300 / 50 = 86$.

\vspace{30pt}

Answer: \colorbox{yellow!100}{C}
\end{solution}

\vspace{60pt}

% Question 31
% Category: Concept | Difficulty: Medium | Question Type: MCQ
\question
A data set consists of 10 identical values. Which of the following must be true?
\begin{tasks}(1)
    \task The mean is 10.
    \task The standard deviation is 0.
    \task The range is 1.
    \task The median is 0.
\end{tasks}

\begin{solution}
\textbf{Logic:} If all values are the same, there is no variation, meaning the spread (Range and SD) is zero.

\vspace{30pt}

Answer: \colorbox{yellow!100}{B}
\end{solution}

\vspace{60pt}

% Question 32
% Category: Exam Level | Difficulty: Medium | Question Type: Numeric Entry
\question
The interquartile range (IQR) of a data set is 15. If the first quartile ($Q1$) is 42, what is the value of the third quartile ($Q3$)?

\begin{solution}
\textbf{Calculation:} $IQR = Q3 - Q1 \Rightarrow 15 = Q3 - 42 \Rightarrow Q3 = 57$.

\vspace{30pt}

Answer: \colorbox{yellow!100}{57}
\end{solution}

\vspace{60pt}

% Question 33
% Category: Practice | Difficulty: Hard | Question Type: MCQ
\question
A set of data has a range of 20. If every value in the set is multiplied by -2, what is the range of the new set?
\begin{tasks}(4)
    \task -40
    \task 20
    \task 40
    \task 60
\end{tasks}

\begin{solution}
\textbf{Logic:} Range is a measure of distance and must always be positive. If all values are multiplied by $k$, the range is multiplied by $|k|$. $20 \times |-2| = 40$.

\vspace{30pt}

Answer: \colorbox{yellow!100}{C}
\end{solution}

\vspace{60pt}

% Question 34
% Category: Concept | Difficulty: Medium | Question Type: MCQ
\question
Which measure is least likely to be affected by an outlier?
\begin{tasks}(4)
    \task Mean
    \task Median
    \task Range
    \task Standard Deviation
\end{tasks}

\begin{solution}
\textbf{Logic:} The median is resistant to outliers because it only depends on the rank/position of values, not their actual numerical weight.

\vspace{30pt}

Answer: \colorbox{yellow!100}{B}
\end{solution}

\vspace{60pt}

% Question 35
% Category: Exam Level | Difficulty: Medium | Question Type: MCQ
\question
In a box plot, what does the rightmost whisker end represent?
\begin{tasks}(4)
    \task Mean
    \task Median
    \task Maximum
    \task Third Quartile
\end{tasks}

\begin{solution}
\textbf{Logic:} The whiskers in a standard SAT box plot extend from the minimum to the maximum values.

\vspace{30pt}

Answer: \colorbox{yellow!100}{C}
\end{solution}

\vspace{60pt}

% Question 36
% Category: Practice | Difficulty: Easy | Question Type: MCQ
\question
A student scored 85, 90, and 88 on three tests. What score does he need on the fourth test to have a mean of 90?
\begin{tasks}(4)
    \task 90
    \task 93
    \task 95
    \task 97
\end{tasks}

\begin{solution}
\textbf{Calculation:} Total needed for 4 tests = $4 \times 90 = 360$. Current total = $85 + 90 + 88 = 263$. $360 - 263 = 97$.

\vspace{30pt}

Answer: \colorbox{yellow!100}{D}
\end{solution}

\vspace{60pt}

% Question 37
% Category: Exam Level | Difficulty: Hard | Question Type: MCQ
\question
Data Set A: {1, 2, 3, 4, 5}. Data Set B: {1, 1, 3, 5, 5}. Which is true?
\begin{tasks}(1)
    \task Mean A $>$ Mean B
    \task Median A $>$ Median B
    \task SD A $>$ SD B
    \task SD B $>$ SD A
\end{tasks}

\begin{solution}
\textbf{Logic:} Both have the same mean (3) and median (3). Set B has more values at the extremes (1s and 5s) compared to A, so Set B has a larger standard deviation.

\vspace{30pt}

Answer: \colorbox{yellow!100}{D}
\end{solution}

\vspace{60pt}

% Question 38
% Category: Practice | Difficulty: Medium | Question Type: Numeric Entry
\question
A data set has 11 values. If the median is 25, how many values must be less than or equal to 25?

\begin{solution}
\textbf{Logic:} In an ordered set of 11, the 6th value is the median. Thus, the 1st through 6th values are all $\le$ the median.

\vspace{30pt}

Answer: \colorbox{yellow!100}{6}
\end{solution}

\vspace{60pt}

% Question 39
% Category: Concept | Difficulty: Medium | Question Type: MCQ
\question
If the standard deviation of a set is 5, what is the variance?
\begin{tasks}(4)
    \task 2.5
    \task 5
    \task 10
    \task 25
\end{tasks}

\begin{solution}
\textbf{Logic:} Variance is the square of the standard deviation. $5^2 = 25$.

\vspace{30pt}

Answer: \colorbox{yellow!100}{D}
\end{solution}

\vspace{60pt}

% Question 40
% Category: Exam Level | Difficulty: Hard | Question Type: MCQ
\question
A frequency table shows 10 students with 0 books, 20 students with 1 book, and $x$ students with 2 books. If the mean is 1, what is $x$?
\begin{tasks}(4)
    \task 5
    \task 10
    \task 15
    \task 20
\end{tasks}

\begin{solution}
\textbf{Calculation:} Mean = Total Books / Total Students.
$1 = \frac{(0\cdot10) + (1\cdot20) + (2\cdot x)}{10 + 20 + x} \Rightarrow 30 + x = 20 + 2x \Rightarrow x = 10$.

\vspace{30pt}

Answer: \colorbox{yellow!100}{B}
\end{solution}

\vspace{60pt}

% Question 41
% Category: Practice | Difficulty: Easy | Question Type: MCQ
\question
What is the mode of {1, 2, 2, 3, 3, 3, 4}?
\begin{tasks}(4)
    \task 1
    \task 2
    \task 3
    \task 4
\end{tasks}

\begin{solution}
\textbf{Logic:} 3 appears most frequently.

\vspace{30pt}

Answer: \colorbox{yellow!100}{C}
\end{solution}

\vspace{60pt}

% Question 42
% Category: Exam Level | Difficulty: Hard | Question Type: MCQ
\question
A set has a mean of 100. If we remove the largest and smallest values, which measure is guaranteed to change?
\begin{tasks}(4)
    \task Mean
    \task Median
    \task Range
    \task Mode
\end{tasks}

\begin{solution}
\textbf{Logic:} Removing the extremes will always shrink the distance between the new max and min, thus changing the range.

\vspace{30pt}

Answer: \colorbox{yellow!100}{C}
\end{solution}

\vspace{60pt}

% Question 43
% Category: Concept | Difficulty: Medium | Question Type: MCQ
\question
In a left-skewed distribution, which is usually true?
\begin{tasks}(1)
    \task mean $<$ median
    \task mean $>$ median
    \task mean = median
    \task SD = 0
\end{tasks}

\begin{solution}
\textbf{Logic:} The tail on the left pulls the mean down below the median.


\vspace{30pt}

Answer: \colorbox{yellow!100}{A}
\end{solution}

\vspace{60pt}

% Question 44
% Category: Practice | Difficulty: Medium | Question Type: MCQ
\question
Find the median of {5, 12, 3, 10, 8}.
\begin{tasks}(4)
    \task 5
    \task 8
    \task 10
    \task 12
\end{tasks}

\begin{solution}
\textbf{Logic:} Ordered: 3, 5, 8, 10, 12. Median = 8.

\vspace{30pt}

Answer: \colorbox{yellow!100}{B}
\end{solution}

\vspace{60pt}

% Question 45
% Category: Exam Level | Difficulty: Hard | Question Type: MCQ
\question
If the sum of 5 numbers is 100, and the median is 20, what is the largest possible value in the set if all numbers are non-negative?
\begin{tasks}(4)
    \task 20
    \task 60
    \task 80
    \task 100
\end{tasks}

\begin{solution}
\textbf{Logic:} To maximize the 5th value, minimize others: 0, 0, 20, 20, $x$. $40 + x = 100 \Rightarrow x = 60$.

\vspace{30pt}

Answer: \colorbox{yellow!100}{B}
\end{solution}

\vspace{60pt}

% Question 46
% Category: Concept | Difficulty: Easy | Question Type: MCQ
\question
Can the mean be greater than the range?
\begin{tasks}(2)
    \task Yes
    \task No
\end{tasks}

\begin{solution}
\textbf{Logic:} Yes. Consider {100, 100, 101}. Mean $\approx 100$, Range = 1.

\vspace{30pt}

Answer: \colorbox{yellow!100}{A}
\end{solution}

\vspace{60pt}

% Question 47
% Category: Practice | Difficulty: Medium | Question Type: Numeric Entry
\question
Median of {1, 2, 3, 4, 5, 100}.

\begin{solution}
\textbf{Calculation:} Even set ($n=6$). Average of 3rd (3) and 4th (4) = 3.5.

\vspace{30pt}

Answer: \colorbox{yellow!100}{3.5}
\end{solution}

\vspace{60pt}

% Question 48
% Category: Exam Level | Difficulty: Hard | Question Type: MCQ
\question
Standard deviation of {5, 5, 5, 5}?
\begin{tasks}(4)
    \task 0
    \task 5
    \task 25
    \task Undefined
\end{tasks}

\begin{solution}
\textbf{Logic:} No variation $\Rightarrow$ SD = 0.

\vspace{30pt}

Answer: \colorbox{yellow!100}{A}
\end{solution}

\vspace{60pt}

% Question 49
% Category: Concept | Difficulty: Medium | Question Type: MCQ
\question
The median of a set is 10. If we add 5 to every number, the new median is?
\begin{tasks}(4)
    \task 10
    \task 15
    \task 50
    \task 5
\end{tasks}

\begin{solution}
\textbf{Logic:} Addition shifts the entire distribution. $10 + 5 = 15$.

\vspace{30pt}

Answer: \colorbox{yellow!100}{B}
\end{solution}

\vspace{60pt}

% Question 50
% Category: Practice | Difficulty: Medium | Question Type: MCQ
\question
If the mean of $x, x+2, x+4$ is 10, what is $x$?
\begin{tasks}(4)
    \task 6
    \task 8
    \task 10
    \task 12
\end{tasks}

\begin{solution}
\textbf{Calculation:} $\frac{3x + 6}{3} = 10 \Rightarrow x + 2 = 10 \Rightarrow x = 8$.

\vspace{30pt}

Answer: \colorbox{yellow!100}{B}
\end{solution}

% =============================================================
% Question End
% =============================================================

\end{questions}
\begin{center}
\rule{.5\textwidth}{1pt}
\end{center}
\end{document}
